\section{Related Work}\label{sec:related}

\subsection{Indic-philosophical architectures in computational settings}

The literature that takes Pratyabhij\~n\=a or adjacent Kashmir \'Saiva architectures as a serious source for AI design is small. \citet{Lawrence1996TantricArgumentPEW} provides the clearest peer-reviewed entry point into the philosophical operators of the Pratyabhij\~n\=a tradition --- in particular into the \emph{discursive} (rather than purely metaphysical) reading of \iast{vimar\'sa} as the moment a speech-act recognises itself as authored from the same ground that contemplates it. \citet{abhinavagupta-isvarapratyabhijna} is the canonical primary source for the \iast{\'sakti}-cascade taxonomy itself, and \citet{pratyabhijnahrdayam-singh} provides Singh's twentieth-century English-translated condensation that has shaped most subsequent secondary literature. The more recent \citet{Waldron2003BuddhistUnconscious} establishes a comparative bridge to the \iast{\=alayavij\~n\=ana} (``store-consciousness'') tradition that PCE uses as the conceptual basis for its consolidation cycle (Appendix~\ref{app:consolidation}).

Our work differs from prior philosophical readings in three load-bearing respects. First, we ground the cascade in concrete operators on a real LLM substrate rather than in a connectionist or theoretical sketch: every \iast{\'sakti} corresponds to a specific computation visible in the Phase 7 traces (\S\ref{sec:engine}). Second, we make the \iast{vimar\'sa} loop a falsifiable trigger under a computable signal (the cascade's $\Delta F$ between draft and shadow revision), not a generic recurrent module; whether it fires correctly is something we can measure, and v0.4 finds it does \emph{not} (H8b). Third, we publish all plugin and benchmark artefacts so any third party can replay the contrast. This last point reflects the Pratyaksha (Skt.\ \emph{pratyak\d{s}a}, ``perception'' / ``directly evidenced'') stance --- the witness invariant of the Pratyaksha Context Engineering Harness~\citep{PratyakshaContextEng2025} from which PCE inherits its discipline that every claim must be anchored to externally verifiable artefacts.

\subsection{Active inference, free-energy, and Bayesian Model Reduction}

The active-inference programme provides the formal substrate that PCE's \iast{j\~n\=ana-\'sakti} (Bayesian-Model-Reduction over a candidate posterior) and \iast{vimar\'sa-\'sakti} (free-energy gated revision) are built on. \citet{Friston2010FreeEnergyUnifiedBrainTheory} introduced the unified free-energy formulation that is now standard in the literature; \citet{Friston2013LifeAsWeKnowItJRSAInterface} situates the same formalism in self-organising systems with Markov-blanket partitions. \citet{Friston2017ActiveInferenceCuriosity} explicitly link the resulting BMR step to insight: a sudden re-organisation of the posterior is the formal correlate of an ``aha'' moment, which is precisely the dynamic PCE intends to instantiate.

\citet{FristonPenny2011} and \citet{Friston2018BayesianModelReductionArxiv} introduced and generalised BMR as a method for posterior switching under nested models, recovering sparse posteriors from rich ones via parameter pruning; \citet{FristonEtAl2016BayesianModelReductionEmpiricalBayesDCM} provides a peer-reviewed empirical-Bayes statement that PCE's \iast{j\~n\=ana} step references for its candidate-posterior reduction. \citet{Heins2022pymdp} (and the recent \texttt{pymdp v1.0} release notes \citep{pymdpRelease100}) provide the discrete-state-space implementation that the v0.5 ladder will use to swap PCE's best-of-K composite-score head for an actual variational free-energy minimisation. \citet{DiPaoloEtAl2024OSFActiveInferenceSchool} has begun making the bridge from active-inference theory to LLM tooling explicit; PCE is one concrete instance of that bridge.

We use BMR as the formal substrate of PCE's \iast{j\~n\=ana} operator (posterior over candidates) and as the conceptual frame for the SWS/REM consolidation cycle (Appendix~\ref{app:consolidation}). PCE's \iast{vimar\'sa} layer is a free-energy-gated revision step that fires when the cascade's $\Delta F$ exceeds a threshold; an honest accounting of what ``free energy'' means in the v0.4 implementation is given in \S\ref{sec:active_inference_background}.

\subsection{LLM-as-judge, scorer-vs-judge calibration, and evaluation methodology}

PCE is built on a two-tier evaluation stack: a small Haiku model serves as the cascade's per-step composite-score head, and a larger Sonnet-4 model serves as a downstream judge against which the scorer is calibrated. The wider literature on this configuration informs our H9 design and our reading of the H9 result. \citet{ZhengEtAl2023JudgingLLMAsJudge} is the canonical reference for the LLM-as-judge protocol and its known failure modes (position bias, verbosity bias, self-enhancement); \citet{LiuEtAl2023GEval} demonstrates that strong LLM judges can serve as reference-free scoring surfaces for open-ended NLG, which is the regime PCE operates in. \citet{SellamEtAl2020BLEURT} establishes the alternative learned-metric lineage (reference-grounded) that any reference-free protocol must be checked against.

Recent work has documented systematic biases that bear directly on PCE's H9 finding (judge $\rho=0$ with scorer). \citet{Braun2025AcquiescenceBiasLLMs} shows that LMs exhibit acquiescence-style answering tendencies that distort rater outputs in ways the scoring surface may absorb differently from the judge. \citet{LauritoEtAl2025PNASAIAIBias} documents an ``AI--AI bias'' in which LMs systematically favour LM-generated text over human-generated text, a property that any judge-of-LM-output protocol must control for. PCE's H9 result --- judge and scorer disagree --- is consistent with this literature; we read it not as a defect of the cascade but as a load-bearing finding for any architecture that uses a small-LM scorer to gate a large-LM commit. \S\ref{sec:discussion} treats this point in detail.

\subsection{Iterative refinement, sample-and-rank, and commit policies}

PCE's commit-policy multiplexer (\texttt{always\_revise} / \texttt{learned\_gate} / \texttt{event\_gated} / \texttt{always\_draft}) sits in a literature converging on the same engineering insight: an LM with a critique step beats an LM without one. \citet{MadaanEtAl2023SelfRefineNeurIPS}'s Self-Refine is the closest mainstream analogue to the cascade's draft$\rightarrow$shadow-revision$\rightarrow$commit pattern; \citet{ShinnEtAl2023Reflexion} extends the analogue to multi-episode reflection buffers; \citet{BaiEtAl2022ConstitutionalAI} establishes the `constitutional' critique-and-revise pattern from which PCE's domain-specific rubric prompts (poetic craft, scientific honesty, etc.) descend. \citet{StiennonEtAl2020LearningToSummarizeFromHumanFeedback} establishes the canonical sample-and-rank lineage that the H6 fairness contrast (\texttt{haiku\_cascade} vs.\ \texttt{haiku\_bare\_2K\_scorer}) is designed to control against: PCE's claim must be that the architecture matters, not merely that doubling the candidate budget at the bare scorer is enough to close the gap.

The literature has not, to our knowledge, isolated the gating question that H8b and H8c address: \emph{when} should a refinement loop revise? Self-Refine and Reflexion default to ``always''; sample-and-rank defaults to ``best-of-K''. PCE's H8c finding --- \texttt{always\_revise} > \texttt{learned\_gate} > \texttt{event\_gated} > \texttt{always\_draft}, with the event-gated $\iast{vimar\'sa}$ trigger as currently defined under-firing --- is, in this context, a small but pointed contribution to the gating sub-literature: at the v0.4 budget, the cheapest correct policy is ``revise unconditionally and let the scorer pick the better of draft and revision''. Whether a properly trained \iast{vimar\'sa} gate can recover the cost saving while keeping the quality is the open question for v0.5 (\S\ref{sec:discussion}).

\subsection{LLM creativity benchmarks}

PCE's four-domain Phase 7 pilot draws on the public creativity-benchmark stack. POEMetric~\citep{Ye2025POEMetric,li2026poemetric} contributes the 6-axis poetry-generation rubric (creativity, lexical diversity, idiosyncrasy, emotional resonance, literary devices, imagery); CreativityPrism~\citep{CreativityPrism2025,hou2025creativityprism} aggregates AUT, TTCT, and other divergent-thinking items under a Quality$\cdot$Novelty$\cdot$Diversity decomposition; \citet{cao2025creval} provides a multi-domain reference-aware text-creativity evaluator; \citet{wang2026creativebench} adds a self-evolving challenges pipeline. \citet{suzgun2023bbh} (BIG-Bench Hard) contributes scientific-reasoning probes; \citet{Tian2024MacGyver} contributes the canonical constraint-satisfaction creativity items. \citet{organisciak2023aut} establishes that LM-based scoring of divergent-thinking outputs can match or exceed human inter-rater agreement on AUT-style items, which underwrites the PCE composite-score head. \citet{Mtenga2024HopfieldCreativity} and \citet{Weber2025SelfOptHopfield} provide the recent Hopfield-network framing of creative-chain dynamics that PCE's deferred Hopfield \iast{\=alayavij\~n\=ana} substrate (Appendix~\ref{app:consolidation}) draws on; \citet{Beaty2015SciRepDMNECN,ChenKenett2025CommunBiolDynamics} establish the empirical neuroscience baseline (DMN/ECN switching during creative production) that motivates a recursive-revision architecture in the first place.

The Wittgenstein aspect-shift task is a custom $n\!=\!20$ stratified slice we construct in \S\ref{sec:benchmarks} on principled grounds, since no public benchmark explicitly tests \emph{simultaneously-coexisting} interpretations of a single line of poetry --- the precise capability the \iast{vimar\'sa} layer is intended to enable.

\subsection{Computational Sanskrit and chandas-aware generation}

The v0.4 release includes a 9-output showcase whose first three items are Sanskrit chandas (anu\d{s}\d{t}ubh, g\=ayatr\={\i}, indravajr\=a) generated under a syllable-count and metre-pattern constraint (\S\ref{sec:showcase}). The validators (\texttt{tools/sanskrit\_chandas.py}) implement minimum-correctness syllabification for both Devan\=agar\={\i} and IAST input; the v0.5 ladder will swap the maintainer-curated Sanskrit references for live cascade outputs once a chandas-aware scorer is wired in. The relevant computational-Sanskrit literature is mature enough to support this: \citet{HellwigEtAl2023VedicSanskritDependencyParsing} and \citet{HellwigBiagetti2025SanskritSembank} establish the corpus-and-annotation backbone, and \citet{NehrdichEtAl2024ByT5Sanskrit} demonstrates a unified ByT5-Sanskrit model that handles segmentation, tagging, and parsing in a single pipeline --- a natural drop-in for the v0.5 chandas-aware scoring head.

\subsection{Plugin architectures for LLMs}

The Cursor and Claude Code ecosystems have converged on a manifest-driven plugin model with MCP tools, slash commands, hooks, agents, and skills. PCE follows the same conventions used by \texttt{attractor-flow}~\citep{AttractorFlow2025} (which couples a real attractor-network engine to MCP), so that the PCE plugin can be installed and audited without any framework-specific shim. \S\ref{sec:substrate} documents the v0.4 portability work: a single source-of-truth plugin declaration mirrored to both Cursor's \texttt{.cursor-plugin/} and Claude Code's \texttt{.claude-plugin/} manifest format, plus a standalone \texttt{python -m pce} CLI wired into \texttt{pyproject.toml}'s \texttt{[project.scripts]} so the substrate can be exercised without any IDE plugin runtime at all.
